% page 365 of the facsimile (image p-364 of the scan)
% block 165.1 x 242.0 mm ; left margin 36.9 mm
\pgc{15.240mm}{0.000mm}{
\l{\cel{53}357}
\l{}
\l{\sou{margrito}. (\sou{bot.}) Nomo di krizantemo kun kolori diversa, qua}
\l{\cel{3}divenis duopla per kultivo - de la astero di Chinia. -}
\l{\cel{3}L.\cel{1}\sou{chrysantemum leucanthemum}. Margriteto : L.\cel{1}\sou{bellis}).(EF.}
\l{}
\l{\sou{mariabalno}. Aquo varma, en vazo, aden qua onu plunjas e mantenas}
\l{\cel{3}korpo quan onu volas varmigar. - DFIS.}
\l{}
\l{\sou{mariajar}. (\sou{trans.}) I. Unionar legitime (viro e muliero), cele-}
\l{\cel{3}brante ta uniono ye la nomo di la lego o di religio. -}
\l{\cel{3}II. - Rezolvigar la mariajo. - EFI.}
\l{}
\l{\sou{marinar}. (\sou{netrans}.) (\sou{aludante karno, fisho)}. Esar en vinagro od}
\l{\cel{3}en vino kun herbi aromatoza, spici, dum kelka tempo ante ko-}
\l{\cel{3}queso. - DEFIRS.}
\l{}
\l{\sou{marioneto}. I. Figurino qua reprezentas persono (viro o muliero)}
\l{\cel{3}quan onu igas agar, marchar, gestifar, e c. per la fili de}
\l{\cel{3}qui lu pendas. - II. (\sou{metaf.}) Persono quan altru direktas,}
\l{\cel{3}igas movar, segun sua arbitrio. - DEFIRS.}
\l{}
\l{\sou{marjino.}\cel{2}Spaco sen-texta, lakuno, cirkum paginedo de texto}
\l{\cel{3}skribura od imprimura. - DEFIS.}
\l{}
\l{\sou{marko}. I. Traco lasita da kozo e per qua onu povas +rekonocar lu.}
\l{\cel{3}II. Traco facita sur kozo por igar lu +rekonocebla. - III.}
\l{\cel{3}Marko obliterala da la posto,qua indikas sur la letri la loko}
\l{\cel{3}e la dio di la departo e di la arivo. - IV. (+timbro).}
\l{\cel{3}Flug-sigluro quan onu glutinas sur letro (por afrankar lu),}
\l{\cel{3}sur afisho, sur quitigo-akto, e c. (por pagar la imposto-}
\l{\cel{3}pekunio) e qua indikas sua preco. - DEFIRS.}
\l{}
\l{\sou{markezo.}\cel{3}I. (\sou{olim)}\cel{2}Chefo di domeno an la frontiero di stato.}
\l{\cel{3}- II. En la hierarkio de la tituli di nobeleso : ta qua venas}
\l{\cel{3}pos la duko ed avan la komto. (Ta vorto esas ligita a la nomo}
\l{\cel{3}di la persono per la prepoziciono \textquotedbl{}de\textquotedbl{}. Ex.: markezo de Beau-}
\l{\cel{3}front). - DEFIRS.}
\l{}
\l{\sou{markotar}. (\sou{trans.)}\cel{3}Enterigar an-sub la surfaco di sulo brancho}
\l{\cel{3}di la planto - brancho quan onu separas de la stipo precipua}
\l{\cel{3}kande lu radifikabos (separo per trancho). - FI.}
\l{}
\l{\sou{marmelado.}\cel{1}Disho ek frukti qui aplastesis e cesis havar sua}
\l{\cel{3}formo normala pro koqueso kun sukro. Anke : la stando di}
\l{\cel{3}karno, di fisho, e c. hiperkoquita. (\sou{anke metaf.)}. DEFIRS.}
\l{}
\l{\sou{marmito.}\cel{2}Implemento koquala - vazo terakota, o metala, kun}
\l{\cel{3}o sen pedo o pedi, en qua onu koquas dishi.\cel{2}- FIS.}
\l{}
\l{\sou{marmoro.}\cel{2}Petro kalkoza blanka o koloroza, veinoza, kun grano}
\l{\cel{3}tenua, polisebla ad-bele, uzata en arkitekturo ed en la}
\l{\cel{3}statuarto. - DEFIRS.}
\l{}
\l{\sou{marmoto}. (\sou{zool}.)Quadripedo rodera, ek la klaso \textquotedbl{}gliri\textquotedbl{}, amanse-}
\l{\cel{3}bla, hibernanta. - L.\cel{1}\sou{arctomys alpina}. - EFIS.}
}
